\section{Research resource architecture and reuse}
\label{sec:resource}

\subsection{Released objects}

The public release separates source evidence, coded evidence, metric inputs, derived results, manuscript claims, and submission artifacts. The ordinary reproduction path uses committed derived records and does not query the forum. Thread-level records contain source links, primary and direct-support codes, evidence type, public outcome, migration status, evidence strength, and analytical notes. Episode records preserve lifecycle stage, technical construct, ecosystem modifiers, evidence form, consequence, public artifact, counterexample status, and confidence.

Field-level metric files contain the case components underlying every reported summary. Robustness files contain alternative partial-score weights, leave-one-thread-out associations, and denominator alternatives. A source-audit ledger maps manuscript and quotation sources to canonical URLs, context review, approved wording, later-update checks, and prohibited inference. A publication claim ledger maps approved manuscript statements to source or data anchors, denominator, safe wording, and explicit overclaim boundaries.

\subsection{Schemas and community records}

The release includes JSON Schemas for a minimum reproducible troubleshooting question, a governed resolved-knowledge index, device-interface records, and run-event streams. The troubleshooting schema activates physical-state and intervention requirements after partial execution. The knowledge schema represents applicability, root-cause status, validation stage, evidence grade, limitations, maintenance ownership, provenance, and lifecycle status. Seed records demonstrate how public discussions can be converted into bounded reusable entries without removing their provenance.

The device-interface model separates documentation, protocol access, licensing, isolated testing, examples, and maintenance. The event model separates commands, acknowledgements, observations, modeled transitions, interventions, recoveries, and final disposition. These schemas are proposed research infrastructure. Their publication is not evidence that they improve operational outcomes; the repository specifies the prospective measures required to evaluate them.

\subsection{Executable reproduction and correction}

Python code recomputes the headline results and association measures from committed records. Regression tests assert published values, schema validity, source mappings, generated-output synchronization, and fail-closed behavior under selected mutations. Figures and LaTeX tables are generated deterministically. Versioned audits record test counts, coverage, source-ledger size, page counts, font checks, and bundle hashes.

Corrections are represented through explicit record status, supersession, and release history. A claim-affecting change must update the source or coded record, dependent metric, generated output, manuscript claim, test, changelog, and versioned release. This control prevents analytical values from changing independently across data, software, and manuscript layers.

\subsection{Recommended reuse}

The resource supports four immediate forms of reuse.
\begin{enumerate}
  \item Researchers can recode the blind hard-case set and report independent agreement without reading the answer-bearing adjudication key.
  \item Laboratory and platform teams can adopt or adapt the question, knowledge, interface, and event schemas.
  \item Researchers can apply the bounded metric pattern to other specialist technical communities while retaining source-specific definitions.
  \item Prospective studies can use the proposed operational outcomes to test the mechanisms identified in the pilot.
\end{enumerate}

Reuse should preserve the sampling frame, unit, denominator, validation stage, and prohibited inference associated with each result. The resource should not be used to rank vendors, infer market share, estimate installed-base failure rates, or profile contributors.
